% =====================================================================
%  c00.tex —— 标题页、译者说明、目录（Lead 独占）
% =====================================================================
\begin{titlepage}
\centering
\vspace*{2.8cm}

{\Huge\bfseries K 理论十八讲\par}
\vspace{0.45cm}
{\Large\itshape 18 Lectures on K-Theory\par}
\vspace{1.0cm}
{\Large Ioannis P. Zois\par}
\vspace{0.4cm}
{\large arXiv:1008.1346v1 [math.KT]，2010 年 8 月\par}

\vspace{1.9cm}
{\LARGE 中 文 译 注 本\par}
\vspace{0.6cm}
{\normalsize 底本：Ioannis P. Zois, \emph{18 Lectures on K-Theory},
  arXiv:1008.1346v1（137 页）\par}

\vfill
{\small
\begin{tabular}{rl}
底本页数 & 137 页（Letter）\\
结构 & 22 节：拓扑 K 理论 6 讲 + 解析 K 理论 6 讲 \\
     & + 高阶代数 K 理论 6 讲 + 附录 \\
编号公式 & 0 条（原文展示式均不编号）\\
翻译体例 & 逐段对译 + 译者批注（不进入排版）\\
\end{tabular}\par}
\vspace{1.3cm}
{\small 本译本仅供个人学习研究使用；原文版权归原作者所有。\par}
\end{titlepage}

\pagenumbering{roman}

% ---------- 译者说明 ----------
\thispagestyle{empty}
\section*{译者说明}
\addcontentsline{toc}{section}{译者说明}

本译本是 Ioannis P. Zois 的《18 Lectures on K-Theory》（arXiv:1008.1346v1，2010）
的中文译注本。原书骨架来自 1995 年伦敦数学会（LMS）在 Lancaster 大学举办的
K 理论研究生暑期学校笔记，其后经作者在 Cardiff 与 Athens 两次研究生课程上扩充而成，
按每小时一讲、每支六讲的方式组织，覆盖 K 理论的三大分支：

\begin{itemize}[leftmargin=1.6em,itemsep=2pt]
  \item \textbf{拓扑 K 理论}（第 2--7 节）：向量丛、分类空间、Bott 周期性、
        Chern 特征与 Chern 类、Adams 运算及其应用（Hopf 不变量为 1 的映射不存在）；
  \item \textbf{解析 K 理论（$K$-同调）}（第 9--14 节）：Fredholm 算子与指标、
        Toeplitz 算子、斜积与配对、$C^*$-代数的 K 理论与扩张；
  \item \textbf{高阶代数 K 理论}（第 16--21 节）：Grothendieck 群、Whitehead 群与
        Steinberg 群、中心扩张、分类空间与群（上）同调、$+$ 构造与一般 $K$ 群。
\end{itemize}

第 22 节为附录，收录代数拓扑、泛函分析与同调代数中若干备用结论；
第 1、8、15 节分别为三支的总起；pp.4--7 为作者自述的「引言与背景」，含一份参考书单。

体例说明：
\begin{itemize}[leftmargin=1.6em,itemsep=2pt]
  \item 正文为忠实逐段翻译，不增不减；原文笔误照译，仅在批注中指出并给出正确形式。
  \item 原文没有编号公式（实测 0 条），故译本亦无；原文的编号定义/命题/引理/定理/推论/注/断言
        一律手工给号，号数照原文（原文每个 section 内从 1 重新编号）。
  \item 原文的展示式（矩阵、分段函数、交换图）全部按原文逐符号重排；交换图用 Ti\emph{k}Z 重画。
  \item 译者批注写成源码中的 \verb|% [Tutor Note: …]| 注释，不进入排版，另见随源码发布的《译注汇总.md》。
  \item 术语首现处括注英文原词；人名、书名与文献键保留原文。
  \item 参考文献（pp.134--137）按学术惯例保留英文原貌，不译。
\end{itemize}

\vspace{0.6em}
\noindent 页码对应：原文 PDF 页号 $=$ 印刷页码（偏移 0）。正文中每起来自新一页处均有
\verb|% [原文 p.N]| 注释，可据此与原文逐页对照。原文 Letter 纸、译本 A4 纸，成品页数与原文不同。

\clearpage
\tableofcontents
\clearpage
